% caltech-memo: LaTeX preamble (XeLaTeX)
% Loaded verbatim via include-in-header, so NO Pandoc template syntax here.

% --- Body font: TeX Gyre Heros (Helvetica-metric, found by filename via
%     kpathsea so it resolves regardless of the render working directory) ---
\usepackage{fontspec}
\setmainfont{texgyreheros}[
  Extension    = .otf,
  UprightFont  = *-regular,
  ItalicFont   = *-italic,
  BoldFont     = *-bold,
  BoldItalicFont = *-bolditalic ]

% Division/department name uses a serif face per Caltech's official lockup
% standard (Adobe Caslon Pro; Georgia is the sanctioned free alternative).
\newfontfamily\divisionfont{Georgia}

\usepackage{graphicx}
\usepackage{xcolor}
\definecolor{caltechorange}{HTML}{FF6C0C}% PMS 1585c (official Caltech orange)
\definecolor{caltechgray}{HTML}{76777B}% PMS Cool Gray 9 (official neutral)

% --- Memo field block ---
% Small gray sans labels (TO / FROM / SUBJECT / DATE / ...), echoing the
% \lfont labels of the old HSSCorrespondence memo environment.
\newcommand{\memolabelfont}{\fontsize{8}{10}\selectfont\color{caltechgray}}
\newlength{\memolabelwd}   % width of the widest label ("SUBJECT")
\newlength{\memorightwd}   % right column (DATE / E-MAIL / MAIL CODE / EXT)
\newlength{\memoleftwd}    % left column (TO / FROM)
\setlength{\memorightwd}{2.2in}
% One line of the right-column contact stack. \par-terminated (not \\) so
% any subset of E-MAIL / MAIL CODE / EXT composes without a trailing break.
\newcommand{\contactline}[2]{\noindent{\memolabelfont #1}\hfill #2\par}

% --- Page furniture (fancyhdr) ---
% Continuation header (page 2+): subject + date left, "Page N" right, over a
% thin Caltech-orange rule (as in the old HSSCorrespondence memo headings).
\usepackage{fancyhdr}
\newcommand{\contheadleft}{}
\fancypagestyle{firstpagestyle}{%
  \fancyhf{}\renewcommand{\headrulewidth}{0pt}}
\fancypagestyle{continuation}{%
  \fancyhf{}%
  \renewcommand{\headrulewidth}{0.5pt}%
  \renewcommand{\headrule}{{\color{caltechorange}\hrule width\headwidth height\headrulewidth}}%
  \fancyhead[L]{\footnotesize\color{caltechgray}\contheadleft}%
  \fancyhead[R]{\footnotesize\color{caltechgray}Page \thepage}}
\pagestyle{continuation}

% Section headings: modest, memo-scaled (the old class set them \bfseries).
\makeatletter
\renewcommand\section{\@startsection{section}{1}{\z@}%
  {-2.2ex \@plus -1ex \@minus -.2ex}{1ex \@plus .2ex}{\large\bfseries}}
\renewcommand\subsection{\@startsection{subsection}{2}{\z@}%
  {-1.8ex \@plus -1ex \@minus -.2ex}{0.8ex \@plus .2ex}{\normalsize\bfseries}}
\makeatother

% Memos are ragged-right, not justified (matches the letter template).
\usepackage{ragged2e}
\RaggedRight
\setlength{\parskip}{0.8\baselineskip}
\setlength{\parindent}{0pt}

% No hyphenation, matching the clean ragged-right look.
\hyphenpenalty=10000
\exhyphenpenalty=10000
